%\documentclass{Math_Note}
%\setlength{\parindent}{2em}

%\begin{document}

% \setcounter{section}{-1}
\newpage
\section{Multiple Integrals Short Sheet (Calculus)}
\subsection{Development of $1$-dim Reimann Integral-Line Integral}
\subsubsection{Intrduction ($ds$)}
\fbox{\begin{minipage}{0.95\textwidth}
\textcolor{skycyan}{Introduction}\\
$I=\int_{C} f(x,y) \, ds,\ C:r=r(\theta),\ \theta \in [\theta_{1},\theta_{2}]$\\
$C:\left\{\begin{matrix} x=r(\theta)\cos\theta \\ y=r(\theta)\sin\theta \end{matrix}\right. ,\ f(x,y)=f(r\cos\theta,r\sin\theta)$\\
$ds=\sqrt{ x'(\theta)^{2}+y'(\theta)^{2} }d\theta=\sqrt{ r'^{2}+r^{2} }d\theta$\\
$I=\int_{\theta_{1}}^{\theta_{2}} f(r\cos\theta,r\sin\theta)\sqrt{ r'^{2}+r^{2} } \, d\theta$
\end{minipage}}
\subsubsection{I-Line Integral}
\fbox{\begin{minipage}{0.95\textwidth}
\textcolor{skycyan}{Definition}\\
$I=\int_{C} f(x,y) \, ds=\left[\left\{\begin{matrix} x=\varphi(t) \\ y=\psi(t) \end{matrix}\right.\right]=\int_{D_{t}} f(\varphi(t),\psi(t))\sqrt{ \varphi'(t)^{2}+\psi'(t)^{2} } \, dt$
\end{minipage}}
\newline\newline
\fbox{\begin{minipage}{0.95\textwidth}
\textcolor{skycyan}{A mean value property}\\
$\rho(x,y)=f(x,y),\ \int_{C(A,B)} f(x,y) \, ds=M,\ \int_{C(C,D)} f(x,y) \, ds=m$\\
$\frac{\int_{C(C,D)} f(x,y) \, ds}{\int_{C(A,B)} f(x,y) \, ds} = \frac{\int_{C(C,D)} C \, ds}{\int_{C(A,B)} C \, ds} = \frac{S_{CD}}{S_{AB}}$
\begin{figure}[H]
    \centering
    \includegraphics[scale=0.36]{"./Figures/250809_1.png"}
\end{figure}
\end{minipage}}
\subsubsection{II-Line Integral}
\fbox{\begin{minipage}{0.95\textwidth}
\textcolor{skycyan}{Definition}\\
$I=\int_{C} f(x,y) \, dx=[x=\varphi(t),y=\psi(t)]=\int_{D_{t}} f(\varphi(t),\psi(t))\varphi'(t) \, dt$\\
$I=\int_{C} f(x,y) \, dy=[x=\varphi(t),y=\psi(t)]=\int_{D_{t}} f(\varphi(t),\psi(t))\psi'(t) \, dt$\\
$I=\int_{C} P \, dx+Q \, dy=[x=\varphi(t),y=\psi(t)]=\int_{D_{t}} [P\varphi'+Q\psi'] \, dt$\\
\end{minipage}}
\subsubsection{Calculate area using II-Line Integral}
\fbox{\begin{minipage}{0.95\textwidth}
\textcolor{skycyan}{Calculate}\\
$D=-\int_{C} y \, dx=\int_{C} x \, dy=\frac{1}{2}\int_{C} x \, dy - y \, dx$, where $C$ is a simple closed cure with positive direction.
\begin{figure}[H]
    \centering
    \includegraphics[scale=0.36]{"./Figures/250809_2.png"}
\end{figure}
\end{minipage}}
\subsubsection{Conditions for curve integrals to be independent of the path}
\fbox{\begin{minipage}{0.95\textwidth}
\textcolor{skycyan}{Theorem}\\
In simply connected space $D$ and given a path $C:A\to B$, II-line integral $\int_{C} P \, dx+Q \, dy$ is dependent of the integration path $C:A\to B$ 
i.f.f. $\frac{\partial P}{\partial y}=\frac{\partial Q}{\partial x}$.
\end{minipage}}
\newline\newline
\fbox{\begin{minipage}{0.95\textwidth}
\textcolor{skycyan}{Proof idea}\\
Necessity "$\Rightarrow$": \\
By studying the differentiation method (partial differentiation with respect to $x$ and $y$) of $F(x, y) = \int_{(x_0, y_0)}^{(x, y)} P  dx + Q  dy$ and the primitive function of $P  dx + Q  dy$ (regarding it as the total differential of some function) (i.e., $\exists \Phi$ such that $P = \frac{\partial \Phi}{\partial x}$, $Q = \frac{\partial \Phi}{\partial y}$ and $\int_C P  dx + Q  dy = \Phi(x, y) \big|_{(x_A, y_A)}^{(x_B, y_B)}$), 
we can conclude that the necessary and sufficient condition for $\int_C P  dx + Q  dy$ to be path-independent is that $P  dx + Q  dy$ ($= d\Phi(x)$) is an exact differential. The latter implies $\frac{\partial P}{\partial y} = \frac{\partial Q}{\partial x}$ (necessity). \\
Sufficiency "$\Leftarrow$": \\
Fix $y \in [c, d]$, integrate $P = \frac{\partial \Phi}{\partial x}$ on $[x_0, x] \subset D$, then $\Phi(x, y) = \int_{x_0}^x P  dx + \Phi(x_0, y)$.
Also, fix $x = x_0$, integrate $Q = \frac{\partial \Phi}{\partial y}$ on $[y_0, y] \subset D$, then $\Phi(x_0, y) = \int_{y_0}^y Q(x_0, y)  dy + \Phi(x_0, y_0)$.\\
Then $\Phi(x, y) = \int_{x_0}^x P(x, y)  dx + \int_{y_0}^y Q(x_0, y)  dy + \Phi(x_0, y_0)$ or $\Phi(x, y) = \int_{y_0}^y Q(x, y)  dy + \int_{x_0}^x P(x, y_0)  dx + \Phi(x_0, y_0)$, which satisfies $\frac{\partial P}{\partial y} = \frac{\partial Q}{\partial x}$.
\end{minipage}}
\subsection{$2$-dim Multiple Integral}
\subsubsection{Definition}
\fbox{\begin{minipage}{0.95\textwidth}
\textcolor{skycyan}{Definition}\\
$\iint_{P} f(x,y) dP=\lim\limits_{ \lambda=\max_{i}\{P_{i}\} \to 0 } \sum\limits_{i} f(\xi_{i},\eta_{i})\Delta P_{i}=V$
\begin{figure}[H]
    \centering
    \includegraphics[scale=0.32]{"./Figures/250809_3.png"}
\end{figure}
\end{minipage}}
\subsubsection{Calculate: $2$-dim Multiple Integral $\to$ Iterated Integral}
\fbox{\begin{minipage}{0.95\textwidth}
\textcolor{skycyan}{$2$-dim Multiple Integral $\to$ Iterated Integral}\\
\underline{Rectangle $P=[a,b]\times[c,d]$}: \\
If $\iint_P f(x,y) dP$ exists and $I_1=\int_a^b f(x,y) dx$ exists $\forall y \in [c,d]$, then
$\iint_P f(x,y) dP=\int_c^d dy \int_a^b f(x,y) dx\ .$\\
If $\iint_P f(x,y) dP$ exists and $I_2=\int_c^d f(x,y) dy$ exists $\forall x \in [a,b]$, then
$\iint_P f(x,y) dP=\int_a^b dx \int_c^d f(x,y) dy\ .$ \\
$\implies$ If $I_1$ and $I_2$ exist, then $\iint_P f(x,y) dP=\int_a^b dx \int_c^d f(x,y) dy=\int_c^d dy \int_a^b f(x,y) dx$.\\
\underline{Curved trapezoid $P=[a,b]\times[y_0(x),Y(x)]$ or $P=[x_0(y),X(y)]\times[c,d]$}: \\
If $\iint_P f(x,y) dP$ exists and $I_2=\int_{y_0(x)}^{Y(x)} f(x,y) dy$ exists $\forall x \in [a,b]$, then
$\iint_P f(x,y) dP=\int_a^b dx \int_{y_0(x)}^{Y(x)} f(x,y) dy\ .$\\
If $\iint_P f(x,y) dP$ exists and $I_1=\int_{x_0(y)}^{X(y)} f(x,y) dx$ exists $\forall y \in [c,d]$, then
$\iint_P f(x,y) dP=\int_c^d dy \int_{x_0(y)}^{X(y)} f(x,y) dx\ .$ \\
$\implies$ If $I_1$ and $I_2$ exist, then $\iint_P f(x,y) dP=\int_a^b dx \int_{y_0(x)}^{Y(x)} f(x,y) dy=\int_{y_0(x)}^{Y(x)} dy \int_a^b f(x,y) dx$.
\begin{figure}[H]
    \centering
    \includegraphics[scale=0.32]{"./Figures/250815_1.png"}
\end{figure}
\end{minipage}}
\subsubsection{Green's Formula: $2$-dim Multiple Integral $\longleftrightarrow$ II-Line Integral}
\fbox{\begin{minipage}{0.95\textwidth}
\textcolor{skycyan}{Green's Formula}\\
$$\iint_{D} (\frac{\partial P}{\partial x}-\frac{\partial Q}{\partial y}) \, dx dy=\oint_{C=\partial D} P \, dy+Q \, dx$$
or 
$$\iint_{D} (\frac{\partial P}{\partial x}+\frac{\partial Q}{\partial y}) \, dx dy=\oint_{C=\partial D} P \, dy-Q \, dx$$
where $D$ is a simply connected domain in $\mathbb{R}^2$ with boundary $C=\partial D$ and $P, Q$ are continuously differentiable functions on $D$.
\end{minipage}}
\newline\newline
\fbox{\begin{minipage}{0.95\textwidth}
\textcolor{skycyan}{Proof idea}\\
$P, Q, \frac{\partial P}{\partial x}, \frac{\partial Q}{\partial y}$ continuously on $D$, $\partial D \triangleq C$.\\
$\iint_{D} \frac{\partial P}{\partial x} dx dy = \int_c^d dy \int_{x_0(y)}^{X(y)} \frac{\partial P}{\partial x} dx = \int_c^d [P(X(y), y) - P(x_0(y), y)] dy = \oint_{\partial D} P(x,y) dy$\\
$\iint_{D} \frac{\partial Q}{\partial y} dx dy = \int_a^b dx \int_{y_0(x)}^{Y(x)} \frac{\partial Q}{\partial y} dy = \int_a^b [Q(x, Y(x)) - Q(x, y_0(x))] dx = - \oint_{\partial D} Q(x,y) dx$
\begin{figure}[H]
    \centering
    \includegraphics[scale=0.36]{"./Figures/250815_2.png"}
\end{figure}
\end{minipage}}
\subsubsection{Variable Substitution in Double Integrals}
\fbox{\begin{minipage}{0.95\textwidth}
\textcolor{skycyan}{Intro}\\
$D=\oint_{S=\partial D} x dy=\left[\begin{cases} x=x(\xi,\eta)\\ y=y(\xi,\eta) \end{cases}\right]=\oint_{\Sigma=\partial\Delta} x(\xi,\eta)\frac{\partial y}{\partial \xi} d\xi + x(\xi,\eta)\frac{\partial y}{\partial \eta} d\eta=\iint_{\Delta} \frac{D(x,y)}{D(\xi,\eta)} d\xi d\eta$
\end{minipage}}
\newline\newline
\fbox{\begin{minipage}{0.95\textwidth}
\textcolor{skycyan}{Analysis}\\
As we could see, the essential of $\begin{cases} x=x(\xi,\eta)\\ y=y(\xi,\eta) \end{cases}$ actually is coordinate transform between the $x,y$-plane and the $\xi,\eta$-plane.\\
$D=\oint_S x dy=\left[\begin{cases} x=x(\xi,\eta)\\ y=y(\xi,\eta) \end{cases},\begin{cases} \xi=\xi(t)\\ \eta=\eta(t) \end{cases} \right]=\int_{\alpha}^{\beta}x(\xi(t),\eta(t))(\frac{\partial y}{\partial\xi}\xi'(t)+\frac{\partial y}{\partial\eta}\eta'(t))dt=\oint_{\Sigma}x(\xi,\eta)\frac{\partial y}{\partial\xi}d\xi+x(\xi,\eta)\frac{\partial y}{\partial\eta}d\eta$,
let $x(\xi,\eta)\frac{\partial y}{\partial\xi} \triangleq P(x,y),\ x(\xi,\eta)\frac{\partial y}{\partial\eta} \triangleq Q(x,y)$, then $D=\oint_{\Sigma} P d\xi + Q d\eta=(\frac{\partial P}{\partial\eta}=\frac{\partial Q}{\partial\xi},\ \text{using Greeni's formula})= \iint_{\Delta}(\frac{\partial Q}{\partial \xi}-\frac{\partial P}{\partial \eta}) d\xi d\eta=\iint_{\Delta}\frac{D(x,y)}{D(\xi,\eta)} d\xi d\eta$.
\end{minipage}}
\newline\newline
\fbox{\begin{minipage}{0.95\textwidth}
\textcolor{skycyan}{Figures for aforementioned analysis}\\
\begin{figure}[H]
    \centering
    \includegraphics[scale=0.32]{"./Figures/250815_3.png"}
\end{figure}
\end{minipage}}
\newline\newline
\fbox{\begin{minipage}{0.95\textwidth}
\textcolor{skycyan}{Idea of Analysis}: \\
Using $\begin{cases} x=x(\xi,\eta)=x(\xi(t),\eta(t))\\ y=y(\xi,\eta)=y(\xi(t),\eta(t)) \end{cases}$ (coordinate transform with parameterization), transform the II-line integration $\oint_Sx dy$ (or $-\oint_S y dx$, or $\frac{1}{2}\oint_S x dy - y dx$) on $x,y$-plane 
to the II-line integral $\oint_\Sigma P(\xi,\eta)d\xi+Q(\xi,\eta)d\eta$, then using Green's formula to transform it to double integral $\iint_\Delta(\frac{\partial Q}{\partial \xi}-\frac{\partial P}{\partial \eta})d\xi d\eta$ on $\xi,\eta$-plane, following is how we 
calculate the double integral (often transform it to iterated integral $\int_{D\xi} d\xi \int_{\eta_0(\xi)}^{\eta(\xi)} (\frac{\partial Q}{\partial \xi}-\frac{\partial P}{\partial \eta})d\xi$, or others).\\
/\textcolor{skycyan}{Shortly,} calculate the area of $D$ on $x,y$ plane: coordinate transform with parameterization $\to$ Green's formula $\to$ Double integration calculate/ \\
\textcolor{skycyan}{Comparison}: \\
Recaping the previous $D=\frac 1 2 \oint_S x dy - y dx$, the original calculation method for this II-line integral is that find thier parametric equation $\begin{cases} x=x(t)\\ y=y(t) \end{cases}, t\in[a,b]$ for different (closed) curve $S$, then substitute into the 
calculation.From here, we could see that the essense of II-line integral is definite integral and has chirality.\\
However, "coordinate transform + Green's formula" method: $D=\iint_{\Delta}\frac{D(x,y)}{D(\xi,\eta)} d\xi d\eta$ as long as we master the coordinate transformation of a type of shape equation, we can easily calculate the results. Although there is a parameterization 
in the derivation, it is not reflected in the result (it can be seen from the process that if the parameterization is continued, it will be the result of the II-line integral in $\xi,\eta$-plane). In essence, the area is calculated as the volume with the same value, 
similiar with $\frac 1 2 \oint_S x dy-y dx=\iint_D 1 dxdy$ or $\oint_S x dy=\iint_D 1 dxdy$ (If you unfold the left and right sides under the curved trapezoid, you can find the consistency.), it's just that the physical meanings of the two equations are different.
\begin{figure}[H]
    \centering
    \includegraphics[scale=0.32]{"./Figures/250815_4.png"}
\end{figure}
/\textcolor{skycyan}{Note}: The choice of which method should be based on which is more convenient, parametric equations or coordinate transformations./
\end{minipage}}
\newline
\fbox{\begin{minipage}{0.95\textwidth}
\textcolor{skycyan}{Variable Substitution in Double Integral}\\
$$\iint_D f(x,y) dxdy=\begin{cases} x=x(\xi,\eta)\\ y=y(\xi,\eta) \end{cases}=\iint_{\Delta} f(x(\xi,\eta),y(\xi,\eta))|\frac{D(x,y)}{D(\xi,\eta)}| d\xi d\eta$$
\end{minipage}}
\subsection{Surface Integral}
\subsubsection{Definition and basics}
\fbox{\begin{minipage}{0.95\textwidth}
\textcolor{skycyan}{Definition}\\
$S=\lim\limits_{\lambda\to 0}\sum\limits_i T_i$, where $T_i$ is the projection area of $S_i$ on the tangent plane at point $M_i\in S_i$.
\begin{figure}[H]
    \centering
    \includegraphics[scale=0.36]{"./Figures/250815_5.png"}
\end{figure}
\end{minipage}}
\newline\newline
\fbox{\begin{minipage}{0.95\textwidth}
\textcolor{skycyan}{Calculte Surface Integral}(Omit matrix-vector operations)\\
$$S=\lim\limits_{\lambda\to 0}\sum\limits_i T_i=\iint_{\Delta} \sqrt{A^2+B^2+C^2} du dv=\iint_{\Delta} \sqrt{EG-F^2} du dv$$
where $A=\begin{vmatrix} y_u' & z_u'\\ y_v' & z_v' \end{vmatrix},\ B=\begin{vmatrix} z_u' & x_u'\\ z_v' & x_v' \end{vmatrix},\ C=\begin{vmatrix} x_u' & y_u'\\ x_v' & y_v' \end{vmatrix}$, 
$E=\sum\limits_{x,y,z} x_u'^2,\ F=\sum\limits_{x,y,z} x_u'x_v',\ G=\sum\limits_{x,y,z} x_v'^2$.
\end{minipage}}
\newline\newline
\fbox{\begin{minipage}{0.95\textwidth}
\textcolor{skycyan}{Procedure:} \\
$\forall M_1\in S'$, the coordinates of its projection $M_1'$ on $T'$ are $\begin{cases} \xi=\xi(u,v)\\ \eta=\eta(u,v)\\ \zeta=\zeta(u,v) \end{cases}$, 
then $S_T'=\iint_{\Delta}|\frac{D(\xi,\eta)}{D(u,v)}| du dv$, then the results could appear after a series of matrix-vector operations.
\begin{figure}[H]
    \centering
    \includegraphics[scale=0.36]{"./Figures/250815_6.png"}
\end{figure}
\end{minipage}}
\subsubsection{Surface Integral}
\fbox{\begin{minipage}{0.95\textwidth}
\textcolor{skycyan}{I-Surface Integral}\\
\begin{equation*}
    \begin{split}
      I &= \iint_S f(x,y,z)\, ds \\
        &= [x=x(u,v),y=y(u,v),z=z(u,v), /parameterization/ \\
        &\ \ \ dS=\sqrt{A^2+B^2+C^2}\, du\, dv=\sqrt{EG-F^2}\, du\, dv] \\
        &= \iint_{\Delta} f(x(u,v),y(u,v),z(u,v)) \sqrt{A^2+B^2+C^2}\, du\, dv \\
        &= \iint_{\Delta} f(x(u,v),y(u,v),z(u,v)) \sqrt{EG-F^2}\, du\, dv
    \end{split}
\end{equation*}
\textcolor{softcyan}{Special case}: \\
If $u=x,\ v=y,\ z=z(x,y)$, \\
\begin{equation*}
  \begin{split}
    I &= \iint_{\Delta} f(x,y,z(x,y)) \sqrt{1+\left(\frac{\partial z}{\partial x}\right)^2+\left(\frac{\partial z}{\partial y}\right)^2} \, dx \, dy \\
      &= \iint_{\Delta} f(x,y,z(x,y)) \frac{dx\, dy}{|\cos(v)|}
  \end{split}
\end{equation*}
\end{minipage}}
\newline\newline
\fbox{\begin{minipage}{0.95\textwidth}
\textcolor{skycyan}{II-Surface Integral} \\
\begin{equation*}
  \begin{split}
    I &= \iint_S f(x,y,z)\, dx\, dy = \text{or} \iint_S f(x,y,z)\, dy\, dz = \text{or} \iint_S f(x,y,z)\, dz\, dx \\
      &\text{or} \iint_S P \, dy\, dz + Q \, dz\, dx + R \, dx\, dy
  \end{split}
\end{equation*}
\end{minipage}}
\newline\newline
\fbox{\begin{minipage}{0.95\textwidth}
\textcolor{skycyan}{Connection between I and II-Surface Integrals} \\
\begin{equation*}
  \begin{split}
    \underset{(II)}{\iint_S f(x,y,z)\, dx\, dy} &= \underset{(I)}{\iint_S f(x,y,z) \cos(v) \, ds} \\
    &\text{(Calculate it: parameterize to get double integral)} \\
    &= [x=x(u,v),y=y(u,v),z=z(u,v), /parameterize/ \\
    \cos&(v)= \pm\frac{C}{\sqrt{A^2+B^2+C^2}},\ dS=\sqrt{A^2+B^2+C^2}\, du\, dv] \\
    &= \pm\iint_{\Delta} f(x(u,v),y(u,v),z(u,v))\cdot C \, du\, dv
  \end{split}
\end{equation*}
\end{minipage}}
\newline\newline
or
\newline\newline
\fbox{\begin{minipage}{0.95\textwidth}
\begin{equation*}
  \begin{split}
    \underset{(II)}{\iint_S P \, dy\, dz + Q \, dz\, dx + R \, dx\, dy} &= \underset{(I)}{\iint_S P\cos(\lambda)+Q\cos(\mu)+R\cos(\nu) \, dS} \\
    \text{(Calculate it: }&\text{parameterize to get double integral)} \\
    &= [\begin{cases} x=x(u,v)\\ y=y(u,v)\\ z=z(u,v) \end{cases} /parameterize/ \\
    &\ \ \ \cos(\lambda)=\pm\frac{A}{\sqrt{A^2+B^2+C^2}},\cdots, \\
    &\ \ \ dS=\sqrt{A^2+B^2+C^2}\, du\, dv] \\
    &= \pm\iint_{\Delta} [P\cdot A+Q\cdot B+R\cdot C] \, du\, dv \\
  \end{split}
\end{equation*}
\end{minipage}}
\subsubsection{Calculate Volume using Surface Integral}
\fbox{\begin{minipage}{0.95\textwidth}
\textcolor{skycyan}{Calculate Volume using Surface Integral}\\
\begin{equation*}
  \begin{split}
    V &= \iint_D Z(x,y) \, dx \, dy-\iint_D z_0(x,y) \, dx \, dy \\
    &= \iint_{S_2} z \, dx \, dy+\iint_{S_1} z \, dx \, dy \\
    &= \iint_S z \, dx \, dy\ (S=S_1+S_2+S_3,\ \text{outside}) \\
    or &= \iint_S x \, dy \, dz \\
    or &= \iint_S y \, dz \, dx \\
    or &= \frac{1}{3}\iint_S (x \, dy \, dz + y \, dz \, dx + z \, dx \, dy)
  \end{split}
\end{equation*}
\begin{figure}[H]
    \centering
    \includegraphics[scale=0.48]{"./Figures/250816_1.png"}
\end{figure}
\end{minipage}}
\subsubsection{Stoke's Formula: II-Surface Integral $\longleftrightarrow$ II-Line Integral}
\fbox{\begin{minipage}{0.95\textwidth}
\textcolor{skycyan}{Stoke's Formula $\Leftarrow$ Development of Green's Formula}\\
$\int_{L=\partial S} P \, dx + Q \, dy + R \, dz = \iint_S (\frac{\partial Q}{\partial x} - \frac{\partial P}{\partial y}) \, dx \, dy + (\frac{\partial R}{\partial y} - \frac{\partial Q}{\partial z}) \, dy \, dz + (\frac{\partial P}{\partial z} - \frac{\partial R}{\partial x}) \, dz \, dx$
\end{minipage}}
\newline\newline
\fbox{\begin{minipage}{0.95\textwidth}
\textcolor{skycyan}{Proof idea}\\
Using Green's formula on LHS as parameterizing $[x=x(u,v), y=y(u,v), z=z(u,v)]$, it's easy to prove.
\end{minipage}}
\subsection{$3$-dim Multiple Integral}
\subsubsection{Definition and basics}
\fbox{\begin{minipage}{0.95\textwidth}
\textcolor{skycyan}{Definition}\\
$\iiint_{V} f(x,y,z) \, dV = \lim\limits_{\lambda\to 0} \sum\limits_i f(x_i, y_i, z_i) \Delta V_i$.
\end{minipage}}
\newline\newline
\fbox{\begin{minipage}{0.95\textwidth}
\textcolor{skycyan}{Calculate}\\
\begin{equation*}
  \begin{split}
    \iiint_{V} f(x,y,z) \, dV &= \int_a^b dx \iint_{P_yz} f(x,y,z) \, dy \, dz \\
    & (\iint_{P_yz} f(x,y,z) \, dy \, dz \ \exists) \\
    &= \int_a^b dx \int_c^d dy \int_e^f f(x,y,z) \, dz \\
    & (\int_e^f f(x,y,z) \, dz \ \exists) \\
    \text{also} &= \iint_{P_xy} \, dx \, dy \int_e^f f(x,y,z) \, dz \\
  \end{split}
\end{equation*}
\begin{equation*}
  \begin{split}
    \iiint_{V} f(x,y,z) \, dV &= \int_{x_0}^{X} dx \iint_{P_yz} f(x,y,z) \, dy \, dz \\
    & (\iint_{P_yz} f(x,y,z) \, dy \, dz \ \exists) \\
    &= \int_{x_0}^{X} \, dx \int_{y_0}^{Y(x)} \, dy \int_{z_0(x,y)}^{Z_(x,y)} f(x,y,z) \, dz \\
    & (\int_{z_0(x,y)}^{Z_(x,y)} f(x,y,z) \, dz \ \exists) \\
    \text{also} &= \iint_{P_xy} \, dx \, dy \int_{z_0(x,y)}^{Z_(x,y)} f(x,y,z) \, dz \\   
  \end{split}
\end{equation*}
\end{minipage}}
\subsubsection{Gauss's Formual $3$-dim Multiple Integral $\longleftrightarrow$ II-Surface Integral}
\fbox{\begin{minipage}{0.95\textwidth}
\textcolor{skycyan}{Gauss's Formula: $\Leftarrow$ Dimension-increasing version of Stokes' formula}\\
\begin{equation*}
  \begin{cases}
    \iiint_{V} \frac{\partial R}{\partial z} \, dx \, dy \, dz = \iint_{S=\partial V} R \, dx \, dy \\
    \iiint_{V} \frac{\partial P}{\partial x} \, dy \, dz \, dx = \iint_{S=\partial V} P \, dy \, dz \\
    \iiint_{V} \frac{\partial Q}{\partial y} \, dz \, dx \, dy = \iint_{S=\partial V} Q \, dz \, dx
  \end{cases}
\end{equation*}
$\Sigma : \implies \iiint_{V} \left( \frac{\partial P}{\partial x} + \frac{\partial Q}{\partial y} + \frac{\partial R}{\partial z} \right) \, dx \, dy \, dz = \iint_{S=\partial V} P \, dy \, dz + Q \, dz \, dx + R \, dx \, dy$
\ \ (S is alonge the outside)
\end{minipage}}
\subsubsection{Variable Substitution in $3$-dim Multiple Integral}
\fbox{\begin{minipage}{0.95\textwidth}
\textcolor{skycyan}{Intro}\\
\begin{equation*}
  \begin{split}
    V &= \iint_S z \, dx \, dy \\
      & \text{(coordinate transform)} \\
      &= \iint_{E} z\cdot\frac{D(x,y)}{D(u,v)} \, du \, dv \\
      &= \iint_{E} z\cdot\left[\frac{D(x,y)}{D(\xi,\eta)}\frac{D(\xi,\eta)}{D(u,v)}+\frac{D(x,y)}{D(\eta,\zeta)}\frac{D(\eta,\zeta)}{D(u,v)}+\frac{D(x,y)}{D(\zeta,\xi)}\frac{D(\zeta,\xi)}{D(u,v)}\right] \, du \, dv \\
      & \text{(II-Surface Integral) (using Green's formula)} \\
      &= \iint_{\Sigma} z\frac{D(x,y)}{D(\xi,\eta)} \, d\xi \, d\eta + z\frac{D(x,y)}{D(\eta,\zeta)} \, d\eta \, d\zeta + z\frac{D(x,y)}{D(\zeta,\xi)} \, d\zeta \, d\xi \\
      &= \iiint_{\Delta} \left(\frac{\partial}{\partial \zeta}\left[z\frac{D(x,y)}{D(\xi,\eta)}\right]+\frac{\partial}{\partial \xi}\left[z\frac{D(x,y)}{D(\eta,\zeta)}\right]+\frac{\partial}{\partial \eta}\left[z\frac{D(x,y)}{D(\zeta,\xi)}\right]\right) \, d\xi \, d\eta \, d\zeta \\
      & \text{(rearrange)} \\
      &= \iint_{\Delta} |\frac{D(x,y,z)}{D(\xi,\eta,\zeta)}| \, d\xi \, d\eta \, d\zeta
  \end{split}
\end{equation*}
\end{minipage}}
\newline\newline
\fbox{\begin{minipage}{0.95\textwidth}
\textcolor{skycyan}{Variable Substitution in $3$-dim Multiple Integral}\\
\begin{equation*}
  \begin{split}
    \iiint_D &f(x,y,z) \, dx \, dy \, dz \\
    &= \left[ \begin{cases} x=x(\xi,\eta,\zeta)\\ y=y(\xi,\eta,\zeta)\\ z=z(\xi,\eta,\zeta) \end{cases} /parameterize/ \right] \\
    &= \iiint_{\Delta} f(x(\xi,\eta,\zeta),y(\xi,\eta,\zeta),z(\xi,\eta,\zeta)) \cdot |\frac{D(x,y,z)}{D(\xi,\eta,\zeta)}| \, d\xi \, d\eta \, d\zeta \\
  \end{split}
\end{equation*}
\end{minipage}}
\subsection{$n$-dim Multiple Integral}
\subsubsection{Definition and basics}
\fbox{\begin{minipage}{0.95\textwidth}
\textcolor{skycyan}{Definition}\\
\begin{equation*}
  \int_{I} f(x^1,\cdots,x^n) \, dx^1 \cdots dx^n = \lim\limits_{\lambda(P)\to 0} \sigma(f,P,\xi)
\end{equation*}
where $I$ is an $n$-dimensional interval, $P$ is a partition of $I$, $\lambda(P)$ is the measure of $P$, and $\sigma(f,P,\xi)$ is the Riemann sum.
\end{minipage}}
\newline\newline
\fbox{\begin{minipage}{0.95\textwidth}
\textcolor{skycyan}{Calculate}\\
\begin{equation*}
  \begin{split}
    \int_{I} f(x^1,\cdots,x^n&) \, dx^1 \cdots dx^n \\
    &\text{(Special)} \\
    &= \int_{a_1}^{b_1} dx^1 \int_{a_2}^{b_2} dx^2 \cdots \int_{a_n}^{b_n} f(x^1,\cdots,x^n) \, dx^n \\
    &\text{(General)} \\
    &= \int_{x_1}^{X_1} dx^1 \int_{x_2(x_1)}^{X_2(x_1)} dx^2 \cdots \int_{x_n(x_1,\cdots,x_{n-1})}^{X_n(x_1,\cdots,x_{n-1}} f(x^1,\cdots,x^n) \, dx^n
  \end{split}
\end{equation*}
\end{minipage}}
\subsubsection{Variable Substitution in $n$-dim Multiple Integral}
\fbox{\begin{minipage}{0.95\textwidth}
\textcolor{skycyan}{Variable Substitution in $n$-dim Multiple Integral}\\
\begin{equation*}
  \begin{split}
    \int_{I} f(x^1,\cdots,x^n&) \, dx^1 \cdots dx^n \\
    &= \left[ \begin{cases} x^1=x^1(\xi^1,\cdots,\xi^n)\\ x^2=x^2(\xi^1,\cdots,\xi^n)\\ \cdots\\ x^n=x^n(\xi^1,\cdots,\xi^n) \end{cases} /parameterize/ \right] \\
    &= \int_{\Delta} f(x^1(\xi^1,\cdots,\xi^n),\cdots,x^n(\xi^1,\cdots,\xi^n)) |\frac{D(x)}{D(\xi)}| \, d\xi^1 \cdots d\xi^n
  \end{split}
\end{equation*}
where $\frac{D(x)}{D(\xi)}=\frac{D(x^1,\cdots,x^n)}{D(\xi^1,\cdots,\xi^n)}$ is the Jacobian determinant of the transformation.
\end{minipage}}

%\end{document}